\section{Computational experiments for higher multipliers}
\label{sec:other-multipliers}

The experiments for \(m\geq4\) show that a single interval is generally not
extremal.  The explicit certificates in Section~\ref{sec:explicit-m4-m5}
establish
\[
    d_4(\T)\geq\frac{11}{64}>\frac16,
    \qquad
    d_5(\T)\geq\frac6{35}>\frac17.
\]
For \(m=5\), searches with up to six intervals have so far only reproduced a
four-interval splitting of the same density \(6/35\); this is evidence, not an
upper bound.

For two intervals, the observed formula is
\begin{equation}
    M_2(m)=\frac{2(m-2)}{m(m+2)},\qquad m\geq4,
    \label{eq:two-interval-conjecture}
\end{equation}
where the restriction \(m\geq4\) is essential because the proved value for
\(m=3\) is \(1/5\).  Formula~\eqref{eq:two-interval-conjecture} agrees with the
reported values \(M_2(6)=1/6\) and \(M_2(7)=10/63\), and it is the \(q=2\)
case of the equally spaced construction in
Section~\ref{sec:equally-spaced-construction}.  No general two-interval upper
bound proving~\eqref{eq:two-interval-conjecture} is currently included.
